
\documentclass[addpoints]{exam}
\usepackage{amsmath}
\usepackage{amssymb}
\usepackage{color}
\usepackage{booktabs}

% \printanswers

\newcommand{\event}{Statistics \& Probability}
\newcommand{\tournament}{}
\def\type{0}

\setlength\fillinlinelength{\textwidth}
\CorrectChoiceEmphasis{\color{red}\bfseries}
\SolutionEmphasis{\color{red}\bfseries}

\setlength\answerclearance{2pt}
\pagestyle{headandfoot}
\runningheadrule
\runningheader{\event}{\tournament}{\ifnum\type=1 Class Set Test \else Full Name:\hspace{1cm} \fi}
\runningfootrule
\runningfooter{}{Page \thepage\ of \numpages}{}

\begin{document}
\begin{center}
    {\huge Grade 11 Foundations for College Math \\ \event
    \ifprintanswers \space Key
    \else \space \ifnum\type<2 Test \fi \ifnum\type=1 \textbf{(CLASS SET)} \fi
          \ifnum\type=2 Answer Sheet \fi
    \fi \\ \vspace{3mm}\tournament\vspace{1mm}}
\end{center}

\vspace{.25\textheight}

\begin{center}
    First Name: \ifprintanswers
    \underline{\hspace{3.5cm}}\textbf{KEY}\underline{\hspace{5.5cm}}\\\vspace{5mm}
    \else \underline{\hspace{10cm}}\\ \vspace{5mm} \fi
    Last Name: \underline{\hspace{10cm}}\\ \vspace{5mm}
\end{center}

\noindent\textbf{\underline{Directions:}}
\begin{itemize}
    \ifnum\type=1 \item \textbf{This test is a class set.} Do not write on this paper. \fi
    \item Show all work for full marks. Round to 2 decimal places unless stated otherwise.
    \item The test is designed to be completed in 75 minutes.
\end{itemize}
\begin{center}
\textit{For grading use only}\\
\multirowgradetable{1}[pages]
\end{center}

\newpage
\begin{questions}

\section*{Part A --- Multiple Choice (15 marks)}
\vspace{5mm}

\question[1] The median of the data set $\{3, 7, 2, 9, 5, 1, 8\}$ is:
\begin{choices}
    \choice $2$
    \choice $3$
    \correctchoice $5$
    \choice $7$
\end{choices}

\question[1] The mode of $\{4, 4, 5, 6, 6, 6, 7, 7\}$ is:
\begin{choices}
    \choice $4$
    \choice $5$
    \correctchoice $6$
    \choice $7$
\end{choices}

\question[1] For the data $\{10, 20, 30, 40, 50\}$, the mean is:
\begin{choices}
    \choice $25$
    \choice $28$
    \correctchoice $30$
    \choice $35$
\end{choices}

\question[1] A data set has $Q_1 = 12$ and $Q_3 = 28$. The IQR is:
\begin{choices}
    \choice $12$
    \choice $14$
    \correctchoice $16$
    \choice $28$
\end{choices}

\question[1] A value is considered an outlier using the $1.5 \times \text{IQR}$ rule if it is:
\begin{choices}
    \choice More than 1 IQR above $Q_3$
    \correctchoice More than $1.5 \times \text{IQR}$ above $Q_3$ or below $Q_1$
    \choice Less than $Q_1$
    \choice More than 2 standard deviations from the mean
\end{choices}

\question[1] The standard deviation of a data set measures:
\begin{choices}
    \choice The middle value
    \correctchoice The spread of data around the mean
    \choice The most common value
    \choice The difference between max and min
\end{choices}

\question[1] A scatter plot shows that as $x$ increases, $y$ decreases. This is:
\begin{choices}
    \choice No correlation
    \choice Strong positive correlation
    \correctchoice Negative correlation
    \choice Positive correlation
\end{choices}

\question[1] The line of best fit for a scatter plot:
\begin{choices}
    \choice Passes through every point
    \correctchoice Passes as close as possible to all the data points
    \choice Has a slope of 1
    \choice Must pass through the origin
\end{choices}

\question[1] A bag contains 4 red, 3 blue, and 3 green marbles. $P(\text{red})$ is:
\begin{choices}
    \choice $\dfrac{3}{10}$
    \choice $\dfrac{1}{4}$
    \correctchoice $\dfrac{4}{10}$
    \choice $\dfrac{4}{6}$
\end{choices}

\question[1] A die is rolled. $P(\text{not a 6})$ is:
\begin{choices}
    \choice $\dfrac{1}{6}$
    \choice $\dfrac{1}{3}$
    \choice $\dfrac{2}{3}$
    \correctchoice $\dfrac{5}{6}$
\end{choices}

\question[1] If $P(A) = 0.4$ and $P(B) = 0.3$ and $A$ and $B$ are mutually exclusive, then
$P(A \text{ or } B) =$:
\begin{choices}
    \choice $0.12$
    \choice $0.10$
    \correctchoice $0.70$
    \choice $1.40$
\end{choices}

\question[1] Theoretical probability is based on:
\begin{choices}
    \choice Conducting many trials
    \correctchoice All equally likely outcomes
    \choice Past results only
    \choice The law of large numbers
\end{choices}

\question[1] A two-way table shows 30 students: 18 play sports, 12 play music, 6 play both.
$P(\text{sports only})$ is:
\begin{choices}
    \choice $\dfrac{18}{30}$
    \correctchoice $\dfrac{12}{30}$
    \choice $\dfrac{6}{30}$
    \choice $\dfrac{24}{30}$
\end{choices}

\question[1] A dataset of exam scores has mean 74 and standard deviation 6. A score of 86 is:
\begin{choices}
    \choice 1 standard deviation above the mean
    \choice 1.5 standard deviations above the mean
    \correctchoice 2 standard deviations above the mean
    \choice 3 standard deviations above the mean
\end{choices}

\question[1] Which correlation coefficient indicates the \textbf{strongest} linear relationship?
\begin{choices}
    \choice $r = 0.5$
    \choice $r = -0.6$
    \correctchoice $r = -0.92$
    \choice $r = 0.3$
\end{choices}


\section*{Part B --- Short Answer (33 marks)}

\question[6] The test scores of 10 students are:
\[52,\; 65,\; 70,\; 72,\; 74,\; 78,\; 80,\; 83,\; 88,\; 90\]
\begin{parts}
    \part[3] Find the mean, median, and range.
    \part[3] Find $Q_1$, $Q_3$, and IQR. Are there any outliers? (Use the $1.5 \times \text{IQR}$ rule.)
\end{parts}
\begin{solution}[9cm]
\textbf{(a)} Mean $= \dfrac{52+65+70+72+74+78+80+83+88+90}{10} = \dfrac{752}{10} = \mathbf{75.2}$

Median: average of 5th and 6th values $= \dfrac{74+78}{2} = \mathbf{76}$

Range $= 90 - 52 = \mathbf{38}$

\textbf{(b)} Lower half: $52, 65, 70, 72, 74$ $\Rightarrow Q_1 = 70$

Upper half: $78, 80, 83, 88, 90$ $\Rightarrow Q_3 = 83$

$\text{IQR} = 83 - 70 = \mathbf{13}$

Lower fence: $70 - 1.5(13) = 70 - 19.5 = 50.5$

Upper fence: $83 + 1.5(13) = 83 + 19.5 = 102.5$

All values are within $[50.5, 102.5]$, so \textbf{no outliers}.
\end{solution}

\question[5] The scatter plot below is described by the following data:

\begin{center}
\begin{tabular}{@{}lcccccc@{}}
\toprule
$x$ (hours studied) & 1 & 2 & 3 & 4 & 5 & 6 \\
\midrule
$y$ (test score) & 45 & 52 & 61 & 68 & 74 & 80 \\
\bottomrule
\end{tabular}
\end{center}
\begin{parts}
    \part[2] Describe the correlation. Estimate a line of best fit equation (approximate).
    \part[3] The line of best fit is $y = 7.2x + 37.4$. Predict the score for 7 hours of study.
    Predict how many hours are needed to score 90. Comment on the reliability of each prediction.
\end{parts}
\begin{solution}[8cm]
\textbf{(a)} The correlation is \textbf{strong and positive} — as study hours increase, test scores increase.

Approximate line of best fit: $y \approx 7x + 38$ (estimates will vary).

\textbf{(b)} At $x = 7$: $y = 7.2(7) + 37.4 = 50.4 + 37.4 = \mathbf{87.8}$

This is interpolation/slight extrapolation — reasonably reliable since it is just beyond the data range.

For $y = 90$: $90 = 7.2x + 37.4 \Rightarrow 7.2x = 52.6 \Rightarrow x \approx \mathbf{7.3}$ hours

This is just beyond the data range — prediction is reasonable but less certain as we extrapolate further.
\end{solution}

\question[6] A survey of 50 students found the following:
\begin{center}
\begin{tabular}{@{}lccc@{}}
\toprule
 & \textbf{Plays Sports} & \textbf{Does Not Play Sports} & \textbf{Total} \\
\midrule
\textbf{Has a Part-time Job} & 10 & 8 & 18 \\
\textbf{No Part-time Job} & 22 & 10 & 32 \\
\textbf{Total} & 32 & 18 & 50 \\
\bottomrule
\end{tabular}
\end{center}
A student is chosen at random.
\begin{parts}
    \part[2] Find $P(\text{plays sports})$ and $P(\text{has a job})$.
    \part[2] Find $P(\text{plays sports AND has a job})$ and $P(\text{plays sports OR has a job})$.
    \part[2] Find $P(\text{does not play sports AND does not have a job})$.
\end{parts}
\begin{solution}[8cm]
\textbf{(a)} $P(\text{sports}) = \dfrac{32}{50} = \mathbf{0.64}$;\quad $P(\text{job}) = \dfrac{18}{50} = \mathbf{0.36}$

\textbf{(b)} $P(\text{sports AND job}) = \dfrac{10}{50} = \mathbf{0.20}$

$P(\text{sports OR job}) = 0.64 + 0.36 - 0.20 = \mathbf{0.80}$

\textbf{(c)} $P(\text{no sports AND no job}) = \dfrac{10}{50} = \mathbf{0.20}$
\end{solution}

\question[8] The daily high temperatures (°C) recorded over 8 days in March are:
\[5,\; 8,\; 6,\; 11,\; 9,\; 7,\; 13,\; 5\]
\begin{parts}
    \part[2] Find the mean $\bar{x}$.
    \part[4] Calculate the population standard deviation $\sigma$ using the formula
    $\sigma = \sqrt{\dfrac{\sum(x_i - \bar{x})^2}{n}}$. Show all steps.
    \part[2] Interpret what the standard deviation means in context. What percent of the data
    values fall within one standard deviation of the mean?
\end{parts}
\begin{solution}[11cm]
\textbf{(a)} $\bar{x} = \dfrac{5+8+6+11+9+7+13+5}{8} = \dfrac{64}{8} = \mathbf{8}$

\textbf{(b)}
\begin{center}
\begin{tabular}{@{}ccc@{}}
\toprule
$x_i$ & $x_i - \bar{x}$ & $(x_i - \bar{x})^2$ \\
\midrule
5 & $-3$ & 9 \\
8 & 0 & 0 \\
6 & $-2$ & 4 \\
11 & 3 & 9 \\
9 & 1 & 1 \\
7 & $-1$ & 1 \\
13 & 5 & 25 \\
5 & $-3$ & 9 \\
\midrule
Total & & 58 \\
\bottomrule
\end{tabular}
\end{center}
\[\sigma = \sqrt{\dfrac{58}{8}} = \sqrt{7.25} \approx \mathbf{2.69 \text{ °C}}\]

\textbf{(c)} The standard deviation of about 2.7°C means the daily temperatures typically vary by
about 2.7°C from the mean of 8°C.

Within one standard deviation: $8 \pm 2.69 \Rightarrow [5.31, 10.69]$.

Values in this range: $8, 6, 9, 7 \Rightarrow$ 4 out of 8 values $= \mathbf{50\%}$.
\end{solution}

\question[8] A student collects data on the number of hours of TV watched per week ($x$) and
the number of hours spent exercising ($y$) for 6 people:

\begin{center}
\begin{tabular}{@{}lcccccc@{}}
\toprule
TV hours $x$ & 2 & 4 & 6 & 8 & 10 & 12 \\
\midrule
Exercise hours $y$ & 9 & 7 & 6 & 5 & 3 & 2 \\
\bottomrule
\end{tabular}
\end{center}
\begin{parts}
    \part[2] Describe the correlation. Is it positive or negative? Strong or weak?
    \part[3] The line of best fit is $y = -0.69x + 10.2$. Use it to predict the exercise hours
    for someone who watches 9 hours of TV per week and for someone who watches 20 hours.
    \part[3] Identify the slope and $y$-intercept and explain what each means in context.
    Discuss any limitations of the model.
\end{parts}
\begin{solution}[9cm]
\textbf{(a)} The correlation is \textbf{strong and negative} — as TV hours increase, exercise hours decrease.

\textbf{(b)} At $x = 9$: $y = -0.69(9) + 10.2 = -6.21 + 10.2 = \mathbf{3.99 \approx 4.0}$ hours — reliable (within data range).

At $x = 20$: $y = -0.69(20) + 10.2 = -13.8 + 10.2 = \mathbf{-3.6}$ hours — \textbf{not valid}; negative exercise hours are impossible. This shows the danger of extrapolating far beyond the data.

\textbf{(c)} Slope $= -0.69$: for each additional hour of TV watched, exercise time decreases by about 0.69 hours (about 41 minutes).

$y$-intercept $= 10.2$: a person who watches 0 hours of TV is predicted to exercise about 10.2 hours/week.

\textbf{Limitations:} The model is only valid for approximately $x \in [2, 12]$. Correlation does not imply causation — other factors (schedule, health, etc.) affect exercise time. The sample size (6) is very small.
\end{solution}

\end{questions}
\end{document}
